\documentclass[11pt,a4paper]{article}
\input{../../shared/preamble}
\SpeedHeader{Geometry}{1}

\begin{document}

\SpeedTitleBlock{Daily Geometry Drill \#1}{Henry Koduthore}
\SpeedMeta{Coordinate geometry: lines, gradients, distances and polygon areas}{negative reciprocal gradients, perpendicular bisectors, distance formula, perpendicular distance to a line, shoelace area}
\noindent\textit{New toolkit today: Gradients \& Perpendicularity $\cdot$ Distance \& Midpoints $\cdot$ Perpendicular Distance $\cdot$ Shoelace Areas $\cdot$ Collinearity.}

\vspace{2pt}
\noindent\textit{Attempt each question before checking answers. Time yourself. No calculators.}

\section*{Section A \quad Rapid Recognition \hfill{\normalfont\small($\leq$15 s each)}}
% ══════════════════════════════════════════════════════════════════════════════

\noindent\textit{Gradient, intercept, midpoint, distance --- read them straight off. No working shown needed.}

\begin{enumerate}

\item Find the gradient of the line $\;6x - 3y + 15 = 0$.

\item Find the gradient of the line through $(1,2)$ and $(5,14)$.

\item A line of gradient $4$ is perpendicular to a line of gradient $m$. Write down $m$.

\item Find the distance between $(2,1)$ and $(5,5)$.

\item Find the midpoint of the segment joining $(4,8)$ and $(10,2)$.

\item Find the perpendicular distance from the origin to the line $\;3x + 4y - 10 = 0$.

\item A line has $x$-intercept $4$ and $y$-intercept $6$. Find its gradient.

\item Find the area of the triangle with vertices $(0,0)$, $(4,0)$ and $(0,3)$.

\item The points $(1,2)$, $(2,4)$ and $(k,10)$ are collinear. Find $k$.

\item Using the shoelace formula, find the area of the triangle with vertices $(0,0)$, $(3,4)$ and $(6,0)$.

\end{enumerate}

% ══════════════════════════════════════════════════════════════════════════════
\section*{Section B \quad Manipulation Drills \hfill{\normalfont\small($\leq$50 s each)}}
% ══════════════════════════════════════════════════════════════════════════════

\noindent\textit{Build the line, then interrogate it. Extra tool required per question.}

\begin{enumerate}

\item The line $\;2x - 3y + 6 = 0\;$ is perpendicular to which of the following?
\begin{itemize}
\item[A)] $y = -\dfrac{3}{2}x + 1$
\item[B)] $y = \dfrac{3}{2}x - 2$
\item[C)] $y = \dfrac{2}{3}x + 5$
\item[D)] $y = -\dfrac{2}{3}x + 3$
\end{itemize}

\item What is the perpendicular distance from $(2,3)$ to the line $\;5x + 12y - 7 = 0$?
\begin{itemize}
\item[A)] $1$
\item[B)] $2$
\item[C)] $3$
\item[D)] $4$
\end{itemize}

\item Which point lies on the perpendicular bisector of the segment joining $(1,2)$ and $(5,6)$?
\begin{itemize}
\item[A)] $(2,5)$
\item[B)] $(3,3)$
\item[C)] $(4,4)$
\item[D)] $(5,3)$
\end{itemize}

\item Find the area of the quadrilateral with vertices $(0,0)$, $(4,0)$, $(5,3)$, $(1,3)$, in that order.
\begin{itemize}
\item[A)] $10$
\item[B)] $12$
\item[C)] $14$
\item[D)] $16$
\end{itemize}

\item The line through $(1,2)$ and $(4,5)$ crosses the $y$-axis at which value?
\begin{itemize}
\item[A)] $0$
\item[B)] $1$
\item[C)] $2$
\item[D)] $3$
\end{itemize}

\item Find the perpendicular distance from $(0,1)$ to the line $\;6x - 8y + 2 = 0$.

\item Which of the following is NOT parallel to $\;y = 3x + 1$?
\begin{itemize}
\item[A)] $y = 3x + 5$
\item[B)] $6x - 2y = 4$
\item[C)] $3x - y = 2$
\item[D)] $y = 2x + 1$
\end{itemize}

\item The line $y = x + 1$ meets the line through $(0,3)$ perpendicular to it at the point $P$. Find the coordinates of $P$.

\item Which point is closest to the origin?
\begin{itemize}
\item[A)] $(1,3)$
\item[B)] $(2,2)$
\item[C)] $(4,0)$
\item[D)] $(0,5)$
\end{itemize}

\item Using the shoelace formula, find the area of the quadrilateral with vertices $(1,0)$, $(5,2)$, $(4,6)$, $(0,3)$, in that order.

\end{enumerate}

% ══════════════════════════════════════════════════════════════════════════════
\section*{Section C \quad Substitution \& Structure \hfill{\normalfont\small($\leq$90 s each)}}
% ══════════════════════════════════════════════════════════════════════════════

\noindent\textit{Spot the tool first --- the numbers were chosen so that a clean route exists. These are genuinely harder: non-obvious structure, hidden invariants, and competition traps.}

\begin{enumerate}

\item The lines $l_1: y = 6 - 2x$ and $l_2: y = \tfrac12 x + 3$ are perpendicular and intersect at $P$. The three lines $l_1$, $l_2$ and the $x$-axis enclose a triangle $T$. What is the area of $T$? \textit{\small(after TMUA 2017 Paper 1 Q3)}
\begin{itemize}
\item[A)] $\tfrac{81}{5}$
\item[B)] $\tfrac{54}{5}$
\item[C)] $\tfrac{81}{10}$
\item[D)] $\tfrac{27}{2}$
\end{itemize}

\item The perpendicular bisector of the segment joining $A(2,-6)$ and $B(5,4)$ meets the $x$-axis at $P$. What is the $x$-coordinate of $P$? \textit{\small(after TMUA Specimen Paper 1 Q3)}
\begin{itemize}
\item[A)] $\tfrac{41}{6}$
\item[B)] $\tfrac13$
\item[C)] $\tfrac{19}{5}$
\item[D)] $\tfrac16$
\end{itemize}

\item The point $A(2,9)$ is reflected in the line $x+2y=5$ to give the point $A'$. How far is $A'$ from the origin?
\begin{itemize}
\item[A)] $\sqrt{85}$
\item[B)] $\sqrt{10}$
\item[C)] $5$
\item[D)] $\sqrt{185}$
\end{itemize}

\item The lines $3x-4y+2=0$ and $4x-3y-5=0$ cross, dividing the plane into four angular regions. Which line bisects the angle whose region contains the point $(1,1)$?
\begin{itemize}
\item[A)] $x-y=7$
\item[B)] $x+y=7$
\item[C)] $7x+7y=3$
\item[D)] $7x-7y=3$
\end{itemize}

\item The points $A(1,4)$, $B(4,7)$ and $C(10,k)$ are collinear. The line through $C$ perpendicular to $AB$ meets the $y$-axis at $D$. What is the $y$-coordinate of $D$?
\begin{itemize}
\item[A)] $23$
\item[B)] $24$
\item[C)] $25$
\item[D)] $26$
\end{itemize}

\item Triangle $ABC$ has vertices $A(0,0)$, $B(9,3)$ and $C(3,9)$. The point $P$ lies on $AB$ with $AP:PB=1:2$, and $Q$ lies on $AC$ with $AQ:QC=2:1$. What is the area of the quadrilateral $PBCQ$?
\begin{itemize}
\item[A)] $24$
\item[B)] $28$
\item[C)] $27$
\item[D)] $32$
\end{itemize}

\item The lines $y=2x+1$, $y=7-x$ and $y=\tfrac12x-2$ enclose a triangle. What is its area?
\begin{itemize}
\item[A)] $12$
\item[B)] $18$
\item[C)] $24$
\item[D)] $48$
\end{itemize}

\item \begin{minipage}[t]{0.68\linewidth}
\vspace{0pt}
The quadrilateral $Q$ has vertices $(0,0)$, $(4,0)$, $(5,3)$ and $(1,3)$ in that order (a parallelogram). A second quadrilateral $Q'$ is obtained by reflecting $Q$ in the line $y=x$. What is the area of the overlap $Q\cap Q'$?
\begin{itemize}
\item[A)] $6$
\item[B)] $8$
\item[C)] $9$
\item[D)] $12$
\end{itemize}
\end{minipage}\hfill
\begin{minipage}[t]{0.28\linewidth}
\vspace{0pt}
\centering
\begin{tikzpicture}[scale=0.42]
\draw[->] (-0.5,0)--(5.8,0) node[right]{\small $x$};
\draw[->] (0,-0.5)--(0,5.8) node[above]{\small $y$};
\draw[speedaux] (-0.3,-0.3)--(5.5,5.5);
\filldraw[fill=gray!15, draw=darkgray] (0,0)--(4,0)--(5,3)--(1,3)--cycle;
\draw[speedgeom, dashed] (0,0)--(0,4)--(3,5)--(3,1)--cycle;
\node at (3.4,0.8) {\small $Q$}; \node at (1.2,3.9) {\small $Q'$};
\node[below] at (4,0) {\tiny $4$}; \node[right] at (5,3) {\tiny $(5,3)$};
\end{tikzpicture}
\end{minipage}

\end{enumerate}

% ══════════════════════════════════════════════════════════════════════════════
\section*{Section D \quad Challenge \hfill{\normalfont\small(TMUA / SMC / BMO1 difficulty)}}
% ══════════════════════════════════════════════════════════════════════════════

\noindent\textit{Hardest 5: deep insight, hidden structure, or a construction you must discover. No diagram gives away the trick.}

\begin{enumerate}

\item The points $A(1,3)$ and $B(7,5)$ are fixed, and $P$ moves along the line $y=x$. What is the least possible value of $PA+PB$?
\begin{itemize}
\item[A)] $4\sqrt2$
\item[B)] $2\sqrt{34}$
\item[C)] $10$
\item[D)] $2\sqrt{10}$
\end{itemize}

\item The square $MNOP$ has side $10$. Points $R$, $S$, $T$, $U$ lie on $MN$, $NO$, $OP$, $PM$ respectively, with $MR=x$, and $RSTU$ is a rectangle whose sides make $45^\circ$ with the sides of the square. For $0<x<10$, what is the largest value of $x$ for which the area of $RSTU$ is $20$? \textit{\small(after TMUA 2023 Paper 1 Q5)}
\begin{itemize}
\item[A)] $5+\sqrt5$
\item[B)] $5+\sqrt{15}$
\item[C)] $5+2\sqrt5$
\item[D)] $10-\sqrt5$
\end{itemize}

\item What is the area of the region of points $(x,y)$ satisfying both $|x|+2|y|\le 6$ and $y\ge x-2$?
\begin{itemize}
\item[A)] $10$
\item[B)] $18$
\item[C)] $26$
\item[D)] $36$
\end{itemize}

\item How many straight lines through the point $(2,3)$ form, together with the two coordinate axes, a triangle of area $12$?
\begin{itemize}
\item[A)] $1$
\item[B)] $2$
\item[C)] $4$
\item[D)] $3$
\end{itemize}

\item For each $t$ with $0<t<4$, the line $L_t$ passes through $(t,0)$ and $(0,4-t)$. What is the greatest possible distance from the origin to $L_t$?
\begin{itemize}
\item[A)] $1$
\item[B)] $2$
\item[C)] $\sqrt2$
\item[D)] $2\sqrt2$
\end{itemize}

\end{enumerate}

\SpeedClosing{``Shoelace before altitude: area is a one-line cross product, not a construction.''}

\end{document}
